\documentclass[11pt]{article}
\usepackage{amsmath,amsthm,amssymb,mathtools}
\usepackage[margin=1in]{geometry}
\usepackage{hyperref}
\usepackage{booktabs}
\usepackage{enumitem}

\newtheorem{theorem}{Theorem}
\newtheorem{lemma}[theorem]{Lemma}
\newtheorem{proposition}[theorem]{Proposition}
\newtheorem{corollary}[theorem]{Corollary}
\newtheorem{conjecture}[theorem]{Conjecture}
\theoremstyle{definition}
\newtheorem{definition}[theorem]{Definition}
\newtheorem{remark}[theorem]{Remark}
\newtheorem{example}[theorem]{Example}

\DeclareMathOperator{\im}{im}
\DeclareMathOperator{\rank}{rank}
\DeclareMathOperator{\supp}{supp}
\DeclareMathOperator{\dep}{dep}
\newcommand{\Vlam}{V^\lambda}
\newcommand{\Ep}[1]{E^{+}_{#1}}
\newcommand{\Em}[1]{E^{-}_{#1}}
\newcommand{\Pq}{P_{q}}
\newcommand{\Pm}{P_{-1}}
\newcommand{\St}{S_{\mathrm T}}
\newcommand{\Sr}{S_{\mathrm R}}

\title{A prefix deficit law for the off--hook order law:\\
each prefix block consumes exactly one level of sign--kill\\
\large (the family $\lambda=(2,2,1^m)$, uniformly in $m$)}
\author{Clio}
\date{30 May 2026 (prove session)}

\begin{document}
\maketitle

\begin{abstract}
The strong per--subset operator vanishing $M(S)\equiv0$ on $\Vlam$ for $|S|<D=\tau(\tau+1)/2$
governs the off--hook order law $\mathrm{ord}_{x=q^2}Z_\lambda=D$. The companion paper
\emph{``A deficit law for Strong Pillar~1''} (29 May 2026) proved the \emph{tail} half: the
tail chain produces sign--kill depth $\tau-|\St|$, i.e.\ $M_{\mathrm T}(\St)\Vlam\subseteq
\bigcap_{j=1}^{\tau-|\St|}\Em j$ for $|\St|<\tau$. Its closing Remark identified the missing
piece as the \emph{prefix counterpart}, controlling $M_{\mathrm R}(\Sr)$. We supply it. Writing
$M(S)=M_{\mathrm R}(\Sr)\,M_{\mathrm T}(\St)$, the prefix acts on the tail's sign--kill image,
and we prove that \emph{each prefix block consumes exactly one level of sign--kill}, at a price
of one $\Pm$--insertion per surviving level. The mechanism is two elementary facts about the
Hoefsmit seminormal module: a sign--kill subspace at depth $d$ forces the first $d$ letters of a
block to be $\Pm$ (else the block annihilates it), and the remaining block letters --- carrying
generators of index $\ge d+1$ --- commute past the generators $1,\dots,d-1$ (axial gap
$\ge2$) and so preserve depth $d-1$. Iterating across the prefix blocks gives the
\textbf{Prefix Deficit Law}: for $\lambda=(2,2,1^m)$, if $|\St|<\tau$ then
\[
   M(S)\neq0\ \Longrightarrow\ |\Sr|\ \ge\ \binom{\tau-|\St|+1}{2}.
\]
Two corollaries follow uniformly in $m$: (i) on the slice $\St=\varnothing$ the full order--law
lower bound, $M(S)=0$ for every prefix--supported $S$ with $|S|<D$; and (ii)
$\mathrm{ord}_{x=q^2}Z_\lambda\ge\tau$. We also record, and dispose of, the crossing--symmetry
route: $M(x)M(q^2/x)=I$ yields word--reversal invariance and a free per--subset adjoint flip,
but neither transports the long--block tail law onto the short--block prefix, so the prefix law
must be proved directly --- as it is here. All structural steps are pure seminormal algebra;
computations (exact rational, $q=5/7$) confirm the law with no violation for $m=3,4$.
\end{abstract}

\section{Setup}\label{sec:setup}

We use the conventions of the 29~May deficit paper. $H_n(q)$ is the Iwahori--Hecke algebra,
$T_i^2=(q-1)T_i+q$; on the Hoefsmit seminormal module $\Vlam$ set $S_i:=T_i+1$ and
\[
   \Ep i:=\ker(T_i-q)|_{\Vlam}=\im(\Pq^{(i)}),\qquad
   \Em i:=\ker(T_i+1)|_{\Vlam}=\im(\Pm^{(i)}),
\]
where $\Pq^{(i)}=(T_i+1)/(q+1)$ and $\Pm^{(i)}=(q-T_i)/(q+1)$ are the orthogonal
eigenprojectors, $\Pq^{(i)}+\Pm^{(i)}=\mathrm{id}$, $\Em i=\ker\Pq^{(i)}$.

\paragraph{The staircase word and the degenerate monodromy.}
For $\lambda=(2,2,1^m)$ we have $n=m+4$, $N=\binom n2$, $\ell(\lambda)=m+2$,
\[
   \tau:=\tau(\lambda)=m-1,\qquad D:=\tfrac{\tau(\tau+1)}2=\binom m2 .
\]
Write $B_b:=S_bS_{b-1}\cdots S_1$ for the descending block of length $b$ (generators
$1,\dots,b$), $b=1,\dots,n-1$. The full staircase operator is
$\Omega=B_1B_2\cdots B_{n-1}$. For a subset $S\subseteq\{1,\dots,N\}$ of word positions the
per--subset operator is
\[
   M(S):=\prod_{k=1}^{N}A_k,\qquad A_k=\Pm^{(i_k)}\ (k\in S),\quad A_k=\Pq^{(i_k)}\ (k\notin S),
\]
$(i_k)$ the staircase letter sequence. As established (operator reformulation),
$\mathrm{ord}_{x=q^2}Z_\lambda=\min\{|S|:M(S)\neq0\}$, so the order law is the assertion that
this minimum equals $D$. We prove a lower bound on it.

\paragraph{Prefix/tail split.}
Following the deficit paper, the tail is the last two (longest) blocks
$B_{n-2}B_{n-1}$ and the prefix is $B_1B_2\cdots B_{n-3}$. Split $S=\Sr\sqcup\St$ into prefix
and tail positions, and factor
\[
   M(S)=M_{\mathrm R}(\Sr)\,M_{\mathrm T}(\St),\qquad
   M_{\mathrm R}(\Sr)=\!\!\prod_{k\in\text{prefix}}\!\!A_k,\quad
   M_{\mathrm T}(\St)=\!\!\prod_{k\in\text{tail}}\!\!A_k .
\]
Because operators act on the right factor first, $M_{\mathrm T}(\St)$ acts on $\Vlam$
\emph{first} and $M_{\mathrm R}(\Sr)$ \emph{afterwards}. The tail deficit law (Corollary~7 of
the deficit paper) supplies the input to the prefix:
\begin{equation}\label{eq:tail}
   |\St|<\tau\ \Longrightarrow\
   M_{\mathrm T}(\St)\,\Vlam\ \subseteq\ \bigcap_{j=1}^{\delta}\Em j,
   \qquad \delta:=\tau-|\St|\ge1 .
\end{equation}

\paragraph{Depth.}
For a subspace $W\subseteq\Vlam$ put
$\dep(W):=\max\{d\ge0: W\subseteq\bigcap_{j=1}^{d}\Em j\}$ (with $\dep(W)=0$ if
$W\not\subseteq\Em1$). Thus \eqref{eq:tail} reads $\dep\big(M_{\mathrm T}(\St)\Vlam\big)\ge\delta$.

\section{The prefix consumes sign--kill, one level per block}\label{sec:engine}

The whole content is two lemmas about how a single descending block $B_b$ acts, by
left multiplication, on a sign--kill subspace. Recall $B_b=S_bS_{b-1}\cdots S_1$, so when applied
to a vector the rightmost letter $S_1$ acts first, then $S_2$, \dots, then $S_b$; with
per--subset projector choices the block becomes $O_b=P^{(b)}P^{(b-1)}\cdots P^{(1)}$,
$P^{(i)}\in\{\Pq^{(i)},\Pm^{(i)}\}$.

\begin{lemma}[A block at depth $d$ costs $d$ insertions]\label{lem:cost}
Let $W\subseteq\Vlam$ with $\dep(W)=d\ge1$, and let $O_b=P^{(b)}\cdots P^{(1)}$ be a block of
length $b\ge d$ acting on $W$. If $P^{(j)}=\Pq^{(j)}$ for some $j\le d$, then $O_bW=0$.
Consequently $O_bW\neq0$ forces $P^{(1)}=\cdots=P^{(d)}=\Pm$, i.e.\ at least $d$ of the block's
letters carry $\Pm$; and then $P^{(d)}\cdots P^{(1)}W=W$.
\end{lemma}

\begin{proof}
Apply the letters in acting order $P^{(1)},P^{(2)},\dots$. Inductively suppose
$P^{(1)},\dots,P^{(j-1)}$ are all $\Pm$ and the running image is still $W$ (true at $j=1$).
Since $j\le d=\dep(W)$ we have $W\subseteq\Em j=\ker\Pq^{(j)}=\im\Pm^{(j)}$. Thus if
$P^{(j)}=\Pq^{(j)}$ then $\Pq^{(j)}W=0$ and the whole product vanishes; while if
$P^{(j)}=\Pm^{(j)}$ then $\Pm^{(j)}W=W$ (a projector restricted to its image is the identity),
and the running image is unchanged. Hence a nonzero outcome forces $P^{(1)}=\dots=P^{(d)}=\Pm$,
in which case $P^{(d)}\cdots P^{(1)}W=W$.
\end{proof}

\begin{lemma}[A block drops depth by at most one]\label{lem:drop}
Let $W\subseteq\Vlam$ with $\dep(W)=d\ge1$ and let $O_b$ ($b\ge d$) be a block with
$O_bW\neq0$. Then $\dep(O_bW)\ge d-1$.
\end{lemma}

\begin{proof}
By Lemma~\ref{lem:cost} the first $d$ letters are $\Pm$ and leave $W$ unchanged, so
$O_bW=P^{(b)}\cdots P^{(d+1)}\,W$ (the letters of index $d+1,\dots,b$, if any; for $b=d$ there are
none and $O_bW=W$, giving depth $d\ge d-1$). Each remaining letter carries a generator of index
$i\in\{d+1,\dots,b\}$. For every $j\in\{1,\dots,d-1\}$ the axial gap satisfies
$|i-j|\ge(d+1)-(d-1)=2$, so $T_i$ commutes with $T_j$, hence $P^{(i)}$ (a polynomial in $T_i$)
commutes with $\Pm^{(j)}$ and therefore preserves $\Em j=\im\Pm^{(j)}$. Thus each $P^{(i)}$,
$i\ge d+1$, maps $\bigcap_{j=1}^{d-1}\Em j$ into itself. Since
$W\subseteq\bigcap_{j=1}^{d}\Em j\subseteq\bigcap_{j=1}^{d-1}\Em j$, applying
$P^{(d+1)},\dots,P^{(b)}$ in turn keeps the image inside $\bigcap_{j=1}^{d-1}\Em j$. Hence
$O_bW\subseteq\bigcap_{j=1}^{d-1}\Em j$, i.e.\ $\dep(O_bW)\ge d-1$.
\end{proof}

\begin{remark}
Lemma~\ref{lem:drop} is sharp: the layer $\Em d$ is generally lost, because the letter $S_{d+1}$
sits at axial gap $1$ from $S_d$ and does not preserve $\Em d$. Numerically the surviving depth
is exactly $d-1$ for the minimal ($\Pm$ only where forced) block.
\end{remark}

\section{The Prefix Deficit Law}\label{sec:main}

We iterate the two lemmas across the prefix blocks. In acting order the prefix blocks are
$B_{n-3},B_{n-4},\dots,B_1$ (the block $B_{n-3}$, nearest the tail, acts first). Their lengths
are $n-3,n-4,\dots,1$, and block $B_{n-2-k}$ (the $k$-th prefix block, $k=1,\dots,n-3$) carries
generators $1,\dots,n-2-k$.

\begin{theorem}[Prefix Deficit Law]\label{thm:main}
For $\lambda=(2,2,1^m)$ and any subset $S$ with $|\St|<\tau$,
\[
   M(S)\neq0\ \Longrightarrow\ |\Sr|\ \ge\ \binom{\delta+1}{2}=\binom{\tau-|\St|+1}{2}.
\]
Equivalently, $M(S)=0$ whenever $|\St|<\tau$ and
$|\Sr|<\binom{\tau-|\St|+1}{2}$.
\end{theorem}

\begin{proof}
Suppose $M(S)\neq0$; then in particular no prefix prefix-of-the-word annihilates the running
image, so the image $W_0:=M_{\mathrm T}(\St)\Vlam$ is nonzero and survives every prefix block.
By \eqref{eq:tail}, $\dep(W_0)\ge\delta$.

Let $W_k$ denote the (nonzero) image after the first $k$ prefix blocks have acted, $W_0$ as
above. We claim, for $k=1,\dots,\delta$: the $k$-th prefix block $B_{n-2-k}$ forces at least
$\dep(W_{k-1})$ distinct $\Pm$--insertions, and $\dep(W_k)\ge\dep(W_{k-1})-1$; combined with
$\dep(W_0)\ge\delta$ this gives $\dep(W_{k-1})\ge\delta-k+1$.

The only hypothesis Lemmas~\ref{lem:cost}--\ref{lem:drop} require of the $k$-th block is that its
length be at least the entering depth $\dep(W_{k-1})$, so that the generators
$1,\dots,\dep(W_{k-1})$ all occur in it. At step $k\le\delta$ the inductive bound gives entering
depth $\dep(W_{k-1})\le\delta-k+1$, while block $B_{n-2-k}$ has length $n-2-k$. Since
$\delta\le\tau=m-1$ and $n-3=m+1$, we have $n-2\ge\delta+1$, hence
\[
   n-2-k\ \ge\ (\delta+1)-k\ \ge\ \delta-k+1\ \ge\ \dep(W_{k-1}),
\]
so the block contains the needed generators. Lemma~\ref{lem:cost} then forces its first
$\dep(W_{k-1})$ letters to be $\Pm$ (that many distinct insertions), and Lemma~\ref{lem:drop}
gives $\dep(W_k)\ge\dep(W_{k-1})-1$.

The insertions counted in different prefix blocks are distinct (they lie in different blocks),
so summing over $k=1,\dots,\delta$,
\[
   |\Sr|\ \ge\ \sum_{k=1}^{\delta}(\delta-k+1)\ =\ \sum_{j=1}^{\delta} j\ =\ \binom{\delta+1}{2}
   \ =\ \binom{\tau-|\St|+1}{2}.
\]
(After $k=\delta$ blocks the depth bound reaches $0$ and no further insertions are forced; the
remaining prefix blocks $B_{n-3-\delta},\dots,B_1$ may be passed with all $\Pq$.) This is the
claim; the contrapositive is the displayed vanishing statement.
\end{proof}

\noindent The block--length bookkeeping in the proof only uses $\delta\le\tau=m-1$ and
$n-3=m+1$, so there are always $\ge\tau$ prefix blocks of sufficient length to host the demanded
$\Pm$'s; the argument is uniform in $m$.

\begin{corollary}[Order--law lower bound on the prefix slice]\label{cor:slice}
For $\lambda=(2,2,1^m)$, $M(S)=0$ for every $S$ supported in the prefix
($\St=\varnothing$) with $|S|<D$. Equivalently, on the slice $\St=\varnothing$ the first nonzero
$M(S)$ occurs at $|S|=D$.
\end{corollary}

\begin{proof}
Take $\St=\varnothing$, so $\delta=\tau$ and Theorem~\ref{thm:main} gives
$M(S)\neq0\Rightarrow|\Sr|\ge\binom{\tau+1}{2}=D$. Hence $|S|=|\Sr|<D$ forces $M(S)=0$.
\end{proof}

\begin{corollary}[$\mathrm{ord}\ge\tau$]\label{cor:ordertau}
For $\lambda=(2,2,1^m)$, $M(S)=0$ for all $|S|<\tau$; hence
$\mathrm{ord}_{x=q^2}Z_\lambda\ge\tau$.
\end{corollary}

\begin{proof}
If $|S|<\tau$ then $|\St|\le|S|<\tau$, so Theorem~\ref{thm:main} applies with
$\delta=\tau-|\St|$. We have $|\Sr|=|S|-|\St|<\tau-|\St|=\delta\le\binom{\delta+1}{2}$ (as
$\binom{\delta+1}{2}\ge\delta$ for $\delta\ge1$), so the survival bound fails and $M(S)=0$.
\end{proof}

\section{The crossing route, and why the prefix law needed a direct proof}\label{sec:crossing}

The Baxterised monodromy $M(x)=\prod_k\check R_{i_k}(x)$, $\check
R_i(x)=\Pq^{(i)}+r(x)\Pm^{(i)}$, $r(x)=(q^2-x)/(q(x-1))$, satisfies the crossing identity
$M(x)M(q^2/x)=I$ (each factor obeys $\check R_i(x)\check R_i(q^2/x)=I$ because
$r(x)r(q^2/x)=1$). One is tempted to transport the tail deficit law onto the prefix by the
word--reversal that crossing induces. We record why this fails, so the route is not re--tried.

\begin{proposition}\label{prop:crossing}
Let $M_{\mathrm{rev}}(x)$ be the monodromy over the reversed staircase word. Then:
\begin{enumerate}[topsep=0.2em,itemsep=0.1em,label=\textup{(\roman*)}]
\item Crossing implies word--reversal invariance, $M_{\mathrm{fwd}}(x)=M_{\mathrm{rev}}(x)$ on
$\Vlam$ for all $x$; matching coefficients of $r(x)^g$ gives only the per--grade identity
$\sum_{|S|=g}M(S)=\sum_{|S|=g}M_{\mathrm{rev}}(S)$, not a per--subset statement.
\item Independently of crossing, the seminormal generators are self--adjoint for the
contravariant form, so $M(S)^{*}=M_{\mathrm{rev}}(S)$; hence $M(S)=0\iff M_{\mathrm{rev}}(S)=0$
(a free per--subset symmetry).
\item The reversed word is the \emph{ascending} staircase; reversal sends the forward prefix
(the \emph{short} blocks $B_1,\dots,B_{n-3}$) to the \emph{end} of the ascending word, whose
trailing blocks are again short. The tail deficit law's engine is the \emph{long}--block hook
structure $B_{n-1}\Vlam=\Ep{n-1}=\Phi(V^{(2,1^m)})$; the short trailing blocks carry no such
structure. So neither \textup{(i)} nor \textup{(ii)} transports the tail law onto the prefix.
\end{enumerate}
\end{proposition}

\begin{proof}
(i) $M(q^2/x)=M(x)^{-1}=\big(\prod_k\check R_{i_k}(x)\big)^{-1}=\prod_k\check
R_{i_{N+1-k}}(q^2/x)=M_{\mathrm{rev}}(q^2/x)$, so $M_{\mathrm{fwd}}=M_{\mathrm{rev}}$ as
functions of the spectral parameter; the coefficient identity is immediate. (ii) Each
$\Pq^{(i)},\Pm^{(i)}$ is self--adjoint for the form making the $T_i$ self--adjoint, and adjunction
reverses products. (iii) Direct inspection of the reversed word. All three are verified
numerically on $(2,2,1,1,1)$ (crossing residual $<10^{-14}$, reversal residual $<10^{-21}$,
adjoint rank--equality on all tested subsets).
\end{proof}

Thus the prefix law is genuinely a new computation, carried out directly in
\S\ref{sec:engine}--\S\ref{sec:main} through the axial--gap commutation of Lemma~\ref{lem:drop}.

\section{Verification}\label{sec:verify}

All checks reuse the Hoefsmit seminormal machinery; structural facts are exact rational
($q=5/7$), bulk sweeps are exact or high--precision numeric at generic $q$. Scripts in
\texttt{\~{}/projects/scratch/2026-05-30-bethe-bridge/}.

\begin{itemize}[topsep=0.2em,itemsep=0.1em]
\item \emph{Tail input \eqref{eq:tail}}: $\dep\big(M_{\mathrm T}(\varnothing)\Vlam\big)=\tau$
exactly for $m=3,4$ ($\Pq^{(j)}$ annihilates the tail image for $j=1,\dots,\tau$).
\item \emph{Lemma~\ref{lem:drop}}: across the first prefix block $B_{n-3}$ at input depth
$\tau$, the minimal--$\Pm$ survivor has output depth exactly $\tau-1$ (no survivor drops $\ge2$).
Verified $m=3,4$.
\item \emph{Theorem~\ref{thm:main}}: the implication
``$|\St|<\tau$ and $|\Sr|<\binom{\tau-|\St|+1}2\Rightarrow M(S)=0$'' holds with no violation
over all subsets $|S|<D$: $67$ law--covered subsets for $m=3$, $6595$ for $m=4$.
\item \emph{Corollary~\ref{cor:slice}}: $M(S)=0$ for all $56$ prefix--supported subsets with
$|S|<D$ at $m=3$, by \emph{exact} rational arithmetic; the first prefix--supported survivor
occurs at $|S|=D=3$.
\item \emph{Corollary~\ref{cor:ordertau}}: $M(S)=0$ for all $|S|<\tau$, exhaustively for
$m=2,3,4$.
\item \emph{Proposition~\ref{prop:crossing}}: crossing, reversal invariance and the adjoint flip
confirmed numerically on $(2,2,1,1,1)$.
\end{itemize}

\begin{center}
\renewcommand{\arraystretch}{1.15}
\begin{tabular}{c|c|c|c|c|c}
\toprule
$\lambda$ & $m$ & $\tau$ & $D$ & order--law slice ($\St{=}\varnothing$) & $\mathrm{ord}\ge\tau$\\
\midrule
$(2,2,1,1)$     & 2 & 1 & 1 & first survivor $|S|{=}1$ & $\ge1$\\
$(2,2,1,1,1)$   & 3 & 2 & 3 & first survivor $|S|{=}3$ & $\ge2$\\
$(2,2,1,1,1,1)$ & 4 & 3 & 6 & first survivor $|S|{=}6$ & $\ge3$\\
\bottomrule
\end{tabular}
\end{center}

\section{Status}\label{sec:status}

\paragraph{Proved (uniformly in $m$).}
\begin{itemize}[topsep=0.2em,itemsep=0.1em]
\item Lemma~\ref{lem:cost} (a depth--$d$ block forces $d$ insertions) --- elementary
eigenprojector identity.
\item Lemma~\ref{lem:drop} (a surviving block drops depth by at most one) --- the only
mechanism, via the axial--gap $\ge2$ commutation of generators $\ge d+1$ past generators
$\le d-1$.
\item Theorem~\ref{thm:main} (Prefix Deficit Law): for $|\St|<\tau$, survival forces
$|\Sr|\ge\binom{\tau-|\St|+1}{2}$. This is the prefix counterpart requested in the Remark of the
29~May deficit paper.
\item Corollary~\ref{cor:slice}: the full order--law lower bound $\mathrm{ord}\ge D$ on the slice
$\St=\varnothing$ (prefix--supported subsets).
\item Corollary~\ref{cor:ordertau}: $\mathrm{ord}_{x=q^2}Z_\lambda\ge\tau$, all subsets.
\item Proposition~\ref{prop:crossing}: the crossing route gives reversal invariance and a free
adjoint flip, but does not transport the tail law to the prefix.
\end{itemize}

\paragraph{Not closed (the precise residual to the full order law $\mathrm{ord}=D$).}
The Prefix Deficit Law is \emph{exact only on the slice $\St=\varnothing$}. The full
$M(S)\equiv0$ for all $|S|<D$ additionally requires:
\begin{enumerate}[topsep=0.2em,itemsep=0.1em,label=\textup{(\Roman*)}]
\item \emph{The intermediate split.} For $1\le|\St|<\tau$ the bound of
Theorem~\ref{thm:main} gives only $|S|\ge g(|\St|)$ with
$g(t)=t+\binom{\tau-t+1}{2}$, and $g(t)<D$ for $t\ge1$ (e.g.\ $g(1)=D-(\tau-1)$). Computation
shows that the subsets realising $|S|=g(t)<D$ \emph{nonetheless vanish} (no survivor among
optimal--prefix $+$ any--tail placements at $|S|<D$ for $m=3,4$), so the lower bound of
Theorem~\ref{thm:main} is \emph{not tight} for $t\ge1$: there is extra killing, presumably from
the placement of the prefix insertions relative to the sign--kill layers, that the per--block
count does not yet capture.
\item \emph{The regime $|\St|\ge\tau$}, outside the tail deficit law's hypothesis, which seeds
the prefix analysis. This is the same open regime flagged in the 29~May paper.
\end{enumerate}
A complete proof of $\mathrm{ord}=D$ thus needs a refinement that (a) charges the
\emph{placement} of prefix insertions, not merely their count, and (b) treats heavy tails; the
present law is the exact backbone of (I) (the $\St=\varnothing$ corner) and reduces the problem
to these two residual phenomena.
\end{document}
